\section{Inference engines and CPU execution}
\label{sec:inference-engines}

An \textbf{inference engine} connects a model's mathematical operations
to executable kernels and manages the data and requests surrounding
those operations.  It loads model representations, prepares inputs,
allocates KV storage and workspaces, dispatches numerical kernels, and
returns generated tokens.  Serving engines additionally manage queues,
request admission and scheduling.  These responsibilities determine when
computation happens and which allocations coexist, even when two
engines evaluate the same model.

\lcpp{} and vLLM are examples with different interfaces, representation
support and execution paths~\cite{llamacppproject2026,vllm2026}.  Their
names do not uniquely specify a kernel or memory policy: backend,
build options, model format and version all matter.  The background uses
engines to illustrate responsibilities; the experimental chapters must
identify the actual implementations and configurations compared.

\subsection{Requests, sequences and execution batches}

\textbf{Concurrency} counts requests or sequences active within a
stated observation boundary.  An \textbf{execution batch} is the set of
positions submitted together in one model evaluation.  A decode batch
can contain one position from each selected sequence, whereas a prefill
batch can contain many positions from one or several prompts.  Queueing,
chunking and admission limits can leave some concurrent requests outside
a particular execution batch.  Request count, sequence count and number
of evaluated positions should therefore be named separately.

For instance, 32 requests can be active while an engine evaluates only
eight sequences in a decode batch.  Alternatively, one prompt can supply
hundreds of positions to a prefill evaluation.  Continuous batching
allows completed sequences to leave and new work to enter between
iterations, changing both the amount of numerical work and the KV
histories being accessed~\cite{kwon2023pagedattention}.  An execution batch
can increase weight reuse, but its effect on response latency depends
also on waiting and scheduling.

\subsection{Kernels, parallelism and precision}

A \textbf{kernel} implements a numerical operation such as a matrix
product, normalization, embedding lookup or attention.  Dense products
with several activation rows can use matrix-matrix operations; a
single activation row gives a matrix-vector-like case.  Optimized
implementations tile work so that subsets of operands can be reused in
caches or registers.  Fusing operations can reduce intermediate writes,
but may change workspace needs and parallelization opportunities.

CPU execution combines instruction-level parallelism with work shared
across cores.  SIMD instructions process several values per instruction;
some processors provide matrix instructions for supported types.
Threads can divide output elements, matrix tiles or attention work.
Their efficiency depends on sufficient independent work, synchronization
cost, load balance and the locality of their operands.  Adding threads
can improve utilization until a compute resource or memory path
saturates; it can also introduce overhead or remote accesses.  Logical
hardware threads are not equivalent to additional physical cores.

The types used for weights, KV storage, temporary activations and
accumulation need not match.  Quantized weights can coexist with BF16
activations and a different KV type.  Lower storage precision changes
the capacity and traffic models, while the supported instructions and
conversion costs determine its execution consequences.  Numerical
accuracy and generated outputs must be evaluated under the chosen
representation rather than assumed identical after quantization.

For performance analysis, phase timers should be complemented by an
understanding of the dispatched kernels and runtime work.  Tokenization,
sampling, scheduling, allocation and synchronization contribute time
without necessarily appearing in a matrix kernel's operation count.
The distinction becomes particularly relevant for small batches or
short requests, where fixed overhead is a larger fraction of elapsed
time~\cite{na2024llmcpu}.
